\section{Recollection on homotopy kan extensions in Model categories}
\subsection{Projective model structure}\label{projms}

\begin{paragr}\label{paragr:htpykan} Let $\M$ be a cofibrantly generated model
category. Recall that for any small category $I$, we can equip the
category of functors $\fun{I}{\M}$ with the \emph{projective model
structure}, where the weak equivalences and the fibrations are these
natural transformations which are pointwise weak equivalences and
fibrations of $\M$ respectively. Any functor $u \colon I \to J$
between small categories induces by precomposition a functor
  \begin{align*} u^* \colon \fun{J}{\M} &\to \fun{I}{\M}\\ X &\mapsto
X\circ u.
  \end{align*} This functor admits a left adjoint $u_!
\colon\fun{I}{\M} \to \fun{J}{\M} $ given by the pointwise Kan
extension formula:
  \[ u_!(X)(j)=\ilim_{\tr{I}{j}}X\vert_{\tr{I}{j}}.
  \] By construction, $u^*$ preserves the weak equivalences and the
fibrations of the projective model structure. Hence, the adjunction
$u_! \dashv u^*$ is a Quillen adjunction and we get an adjunction at
the level of homotopy categories
  \[ \LL u_! \colon \Loc{\fun{I}{\M}} \leftrightarrows
\Loc{\fun{J}{\M}} \colon u^*.
  \] Here, we kept the same notation $u^*$ at the level of homotopy
categories (instead of, say, $\mathbb{R} u^*$) because the functor
$u^*$ preserves weak equivalences and thus we have a commutative
square
  \[
    \begin{tikzcd} \fun{J}{\M} \ar[d] \ar[r,"u^*"] &
\fun{I}{\M}\ar[d]\\ \Loc{\fun{J}{\M}} \ar[r,"u^*"] &
\Loc{\fun{I}{\M}},
    \end{tikzcd}
  \] where the vertical arrows are the localization functors.
\end{paragr}
\begin{paragr} We denote by $\Term$ the terminal category and make the
canonical identification $\fun{\Term}{\M}=\M$. For a small category
$I$, we denote by $p_I \colon I \to \Term$ the unique such
functor. The functor $p_I^* \colon \M \to \fun{I}{\M}$ is the
``diagonal'' functor that associates to an object $X$ of $\M$, the
constant diagram $I \to \M$ with value $X$. Hence, $p_{I!}$ is nothing
but the colimit functor $\ilim_I \colon \fun{I}{\M} \to \M$, and we
shall use the two notations interchangeably. Consequently, $\LL
p_{I!}$ is nothing but the \emph{homotopy colimit functor}
  \[ \hoilim_I \colon \Loc{\fun{I}{\M}} \to \Loc{\M}.
  \]
\end{paragr}
\begin{paragr} Let $I$ be a small category and $i \in \Ob(I)$ an
object of $I$. This can be seen as a functor $i \colon \Term \to I$,
which induces a ``evaluation at $i$'' functor
  \begin{align*} i^* \colon \fun{I}{\M} &\to \M\\ X &\mapsto X(i).
  \end{align*} It is a standard fact of category theory that the
family of functors
  \[\left(i^*\colon \fun{I}{\M} \to \M\right)_{i\in \Ob(I)}\] is
\emph{conservative}, meaning that a morphism $\alpha \colon X \to Y$
of $\fun{I}{\M}$ is an isomorphism if and only if for each object $i$
of $I$, $i^*(\alpha)=\alpha_i$ is an isomorphism. The same is true at
the level of homotopy categories as we recall in the following result.
\end{paragr}
\begin{proposition} Let $\M$ be a (cofibrantly generated) model
category and $I$ a small category. The family of functors
  \[ \left(i^* \colon \Loc{\fun{I}{\M}}\to \Loc{\M}\right)_{i \in
\Ob(I)}
  \] is conservative.
\end{proposition}
\begin{proof}
  
\end{proof}
\subsection{Homotopy cocontinuous functors}\label{hptycocts}
\begin{paragr} Let $F \colon \M \to \M'$ be a left Quillen functor
between (cofibrantly generated) model categories. For any small
category $I$, we consider the functor induced by post-composition
  \begin{align*} F^I \colon \fun{I}{\M} &\to \fun{I}{\M'}\\ X &\mapsto
F\circ X.
  \end{align*} This functor is also a left Quillen functor with
respect to the projective model structures. Hence, it induces a left
derived functor
  \[ \LL F^I \colon \Loc{\fun{I}{\M}}\to \Loc{\fun{I}{\M'}}.
  \]

  Now, let $u \colon I \to J$ be a functor between small
categories. By construction, we have a commutative square
  \[
    \begin{tikzcd} \fun{J}{\M} \ar[r,"F^J"] \ar[d,"u^*"']&
\fun{J}{\M'} \ar[d,"u^*"]\\ \fun{I}{\M} \ar[r,"F^I"'] & \fun{I}{\M'},
    \end{tikzcd}
  \] and, because $F$ is a left adjoint, the following square is
commutative up to a canonical isomorphism
  \[
    \begin{tikzcd} \fun{I}{\M} \ar[r,"F^I"] \ar[d,"u_!"']&
\fun{I}{\M'} \ar[d,"u_!"]\\ \fun{J}{\M} \ar[r,"F^J"'] & \fun{J}{\M'}.
\ar[from=1-1,to=2-2,phantom,"\simeq" description]
    \end{tikzcd}
  \]
  This fact is usually referred to as ``$F$ preserves Kan
extension''. We recall now the following result which says that the
same holds at the level of homotopy categories for a left Quillen
functor.
\end{paragr}
\begin{proposition} Let $\M$ and $\M'$ be two (cofibrantly generated)
  model categories, $F \colon \M \to \M'$ be a left Quillen functor and
  $u \colon I \to J$ a functor between small categories. Then the
  squares
  \[
    \begin{tikzcd} \Loc{\fun{J}{\M}} \ar[r,"\LL F^J"] \ar[d,"u^*"']&
      \Loc{\fun{J}{\M'}} \ar[d,"u^*"]\\ \Loc{\fun{I}{\M}} \ar[r,"\LL F^I"']
      & \Loc{\fun{I}{\M'}},
    \end{tikzcd}
  \] and
  \[
    \begin{tikzcd} \Loc{ \fun{I}{\M}} \ar[r,"\LL F^I"] \ar[d,"\LL
      u_!"']& \Loc{\fun{I}{\M'}} \ar[d,"\LL u_!"]\\ \Loc{\fun{J}{\M}}
      \ar[r,"\LL F^J"'] & \Loc{\fun{J}{\M'}},
    \end{tikzcd}
  \] are commutative up to a canonical isomorphism.
\end{proposition}
\begin{proof}
\end{proof}
\subsection{Abstract homology of small categories}\label{abstracthom}
\begin{definition}\label{def:abstracthom}
  Let $\M$ be a (cofibrantly generated) model category. For a small
  category $I$ and $X$ an object of $\fun{I}{\M}$, we define the \emph{$\M$\nbd-valued
    homology of $I$ with coefficient $X$} as the object of $\Loc{\M}$
  \[
    \Ho_{\M}(I,X):=\LL p_{I!}X=\hoilim_IX.
  \]
\end{definition}
\begin{remark}
  Note that the use of the word ``homology'' here is arguably slightly
  abusive as we don't make any hypothesis of additivity or stability
  property of any kind on (the localization of) $\M$. 
\end{remark}
\begin{paragr}
 Obviously, for a fixed object small category $I$, the previous
 construction yields a functor
 \[
   \Ho_{\M}(I,-) \colon \Loc{\fun{I}{\M}}\to \Loc{\M}.
 \]
 Let us now explore the functoriality of $\Ho_{\M}(I,X)$ in $I$. Let
  $u \colon I \to J$ be a functor between small categories. We have a
  commutative triangle
  \[
    \begin{tikzcd}[column sep=small]
      I \ar[rr,"u"]\ar[rd,"p_I"'] && J\ar[ld,"p_J"] \\
      &\Term&
    \end{tikzcd}
  \]
  which induces a commutative triangle
  \[
    \begin{tikzcd}[column sep=small]
      \Loc{\fun{I}{\M}}  && \Loc{\fun{J}{\M}}\ar[ll,"u^*"'] \\
      &\Loc{\M}.\ar[lu,"p_I^*"]\ar[ru,"p_J^*"']&
    \end{tikzcd}
  \]
  Hence, by uniqueness of adjoints, this yields a triangle commutative
  up to a canonical isomorphism
  \[
        \begin{tikzcd}[column sep=small]
      \Loc{\fun{I}{\M}} \ar[rr,"\LL u_!"]\ar[rd,"\LL p_{I!}"'] &&
      \Loc{\fun{J}{\M}}\ar[ld,"\LL p_{J!}"] \\
      &\Loc{\M}.&
    \end{tikzcd}
  \]
  Now let $X$ be an object of $\fun{J}{\M}$. Then we have a morphism
  \[
    \LL p_{I!}u^*X \simeq \LL p_{J!} \LL u_! u^*X\to \LL p_{J!}X,
  \]
  where the morphism of the right is the co-unit of the ajdunction
  $\LL u_! \dashv u^*$. On other words, this defines a morphism
  \[
  \Ho_{\M}(u,X)\colon  \Ho_{\M}(I,u^*X)\to \Ho_{\M}(J,X).
  \]
  We then have the following functoriality result.
\end{paragr}
\begin{proposition}\label{prop:functorialityhomology}
  Let $\M$ a (cofibrantly generated) model category, $K$ a small
  category and $X$ an object of $\fun{K}{\M}$. We have a functor
  \[
      \Ho_{\M}(-,X) \colon \tr{\Cat}{K} \to \Loc{\M},
    \]
    whose action on object is 
  \[
    (I,p \colon I \to K) \longmapsto \Ho_{\M}(I,p^*X)
  \]
  and whose action on morphisms is
  \[
    u\colon (I,p)\to (J,q) \longmapsto \Ho_{\M}(u,q^*X)\colon \Ho_{\M}(I,p^*X) \to \Ho_{\M}(J,q^*X).
\]
 
\end{proposition}
\begin{proof}
\end{proof}

\begin{paragr}
  The main case of interest of the previous proposition will be in the
  case that $K=\Term$ is the terminal category. Then $F$ is simply an
  object $X$ of $\M$ and we can make the identification
  $\tr{\Cat}{K}\simeq \Cat$. We speak of the (abstract) homology of a
  small category $I$ with \emph{constant} coefficient $X$ for
  \[
    \Ho(I,p_I^*X).
  \]
  \cref{prop:functorialityhomology} in this case yields the following corollary.
\end{paragr}
\begin{corollary}
  Let $\M$ be a (cofibrantly generated) model category and $X$ an
  object of $\M$. We have a functor
  \[
    \Ho(-,X) \colon \Cat \to \Loc{\M},
  \]
  whose action on objects is
  \[
    I \longmapsto \Ho_{\M}(I,p_I^*X),
  \]
  and whose action on morphisms is
  \[
    u \colon I \to J \longmapsto \Ho_{\M}(u,X) \colon
    \Ho_{\M}(I,p_I^*X) \to \Ho_{\M}(J,p_J^*X).
  \]
\end{corollary}
\subsection{Homology of categories}
\begin{paragr}\label{paragr:homalg}
  For $C$ a small category, we denote by $\LMod{C}$ the category of
  \emph{right $C$\nbd-modules}, that is of functors
  \[
    C \to \Ab
  \]
  from $C$ to the category of abelian groups, and of natural
  transformations between them. By definition, the category of \emph{right}
  $C$\nbd-modules is the category
  \[
    \RMod{C}:=\LMod{C^{\op}}
  \]
  and everything that follows is dualized accordingly. 
  
  We denote by $\Chp(C)$ the category of chain complexes
  of left $C$\nbd-modules in non-negative degree. For $C=\Term$ the
  terminal category, we use the notation $\Chp(\Z)$ instead. 
  We denote by $\Derp{C}$ the \emph{derived
 category of $C$} in non-negative degree, that is the
localization of $\Chp(C)$ with respect to
 quasi-isomorphisms (and we use the notation $\Derp{\Z}$ when
 $C=\Term$ is the terminal category). This localization is controlled by the
  so-called \emph{projective} model structure on $\Chp(C)$, which is
  cofibrantly generated,
  characterized as follows:
  \begin{itemize}
  \item the weak equivalences are the quasi-isomorphisms,
  \item the fibrations are the levelwise epimorphisms in each positive
    degree,
  \item the cofibrations are the levelwise monomorphisms with
    projective cokernel in each non-negative degree.
  \end{itemize}
  In particular, cofibrant objects are the chain complexes (in
  non-negative degree) of projective
  $C$\nbd-modules. The following lemma is useful to
  construct and recognize such objects.
\end{paragr}
\begin{lemma}\label{lemma:freeisproj}
  Let $c$ be an object of $C$. The left $C$\nbd-module $\Z[\homset{C}{c}{-}]$
  \[
    \begin{aligned}
      C &\to \Ab\\
      d &\mapsto \Z[\homset{C}{c}{d}],
    \end{aligned}
  \]
  is a projective object of $\LMod{C}$.
\end{lemma}
\begin{proof}
  This follows straightforwardly from the Yoneda lemma.
\end{proof}

\begin{paragr}
   Notice that for a small category $C$, we have a canonical
   \emph{isomorphism} of categories
  \[
    \Chp(C)\simeq \fun{C}{\Chp(\Z)},
  \]
  and we shall make this identification implictly. Note that this identification is compatible with
  model structures in the sense that the projective model structure
  on $\fun{C}{\Chp(\Z)}$ in the sense of \ref{paragr:htpykan} is the
  same, via the previous isomorphism, as the projective model
  structure on $\Chp(C)$ in the sense of \ref{paragr:homalg}.

  This allows us to use all the results from the sections
  \ref{projms}, \ref{hptycocts} and \ref{abstracthom}. 
\end{paragr}
\begin{definition}
  Let $C$ be a small category and $X$ a left $C$\nbd-module. We define
  the \emph{homology of $C$ with coefficient $X$}, and denote by
  \[
    \Ho(C,X),
  \]
  the $\Chp(\Z)$\nbd-valued homology of $C$ with coefficient $X$
  (\cref{def:abstracthom}).
\end{definition}
%%% Local Variables:
%%% mode: LaTeX
%%% TeX-master: "main"
%%% End:
